% !PRINTED 59
\weberpage{59}

If one replaces $\varphi(u)$ by $\varphi(u)-c$, where $c$ denotes an
arbitrary value, the important theorem follows:

3. \gesperrt{The sum of the argument values at which a
doubly periodic function takes one and the same value $c$
 in the interior of a period parallelogram
is independent of $c$, or under continuous variation
 of $c$ changes at most by jumps of a period.}

4. \gesperrt{From this it follows that there is also no doubly
 periodic function of the first order.}

For such a function would have to take every value $c$ at a single
point $u$. If, then, it is equal to $c'$ at a point $u'$
and equal to $c''$ at a point $u''$ different from it, it would
follow from 3. that $u'\equiv u''$, contrary to the assumption.

By a \gesperrt{doubly periodic function of the second
kind} one understands, following \gesperrt{Hermite}, a function $\psi(u)$ which
behaves with respect to the periods $\omega_1$, $\omega_2$ according to the equations:

\[ \psi(u+\omega_1) = a_1\,\psi(u)\,, \quad \psi(u+\omega_2) = a_2\,\psi(u)\,, \]

where $a_1$, $a_2$ are constants.

5. \gesperrt{For a doubly periodic function of the
second kind also the theorem holds that it becomes zero
and infinite at equally many points of the period
 parallelogram.} If the number of these points is $m$, then here too the
function is said to be of the $m$th order.

The proof is the same as that of theorem 2. But the
extension no longer holds here, that the function $\psi(u)$ takes every
value equally often, because $\psi(u)-c$ is now no longer
 periodic.

6. \gesperrt{A doubly periodic function $\psi(u)$ of the
second kind and of order zero is necessarily a
constant or an exponential function.} For if $\psi(u)$
does not become zero, then

\[ \varphi(u) = \dfrac{\psi'(u)}{\psi(u)} \]

is a doubly periodic function of the first kind which does not
become infinite, and hence, by 1., a constant $a$. From this it follows:

\[ \psi(u) = A\,e^{a\,u}\,, \]

where $A$ is a new constant.
